The first three parts discussed the  AI techniques for constructing advanced LLM Agents that can interact with and act in the real world. These capabilities allow them to do useful things for their users. In this  part we will address the use of tools by LLMs to perform specialized tasks, such as looking up information, or controlling robots.  Agents use tools through learned APIs (Application Programming Interfaces).

Tools can help us solve a major problem of pretrained LLMs, which is that the information that is contained in them is only as up to date as the dataset on which they were trained. Retrieval augmentation techniques allows LLMs to provide up to the minute information.  

Many types of such agents have been created, and in the final chapter of this part we will review important examples, from the medical, financial, and scientific domains. Finally, we will discuss a practical case study in which the use of the technologies comes together, when we will implement an LLM Agent to perform medical diagnosis.
